% Checks cases the fix must not break: test-hypcap without hypcap, whose
% numbers and references the page with hypcap must match.
\documentclass{article}
\usepackage{float}
\usepackage{hyperref}
\usepackage[caption=false]{subfig}
\usepackage{cleveref}
\newfloat{program}{tbp}{lop}
\floatname{program}{Program}
\crefname{program}{program}{programs}
\newtheorem{theorem}{Theorem}
\begin{document}
\section{Results}\label{sec:res}
\begin{figure}[h]
  \centering
  \subfloat[Left half]{\fbox{Left}\label{fig:a}}\quad
  \subfloat[Right half]{\fbox{Right}\label{fig:b}}
  \caption{First figure}\label{fig:one}
\end{figure}
\begin{figure}[h]
  \centering
  \captionsetup{name=Fig.}
  \fbox{Second}
  \caption{Second figure}\label{fig:two}
\end{figure}
\begin{table}[h]
  \centering
  \caption{First table}\label{tab:one}
  \subfloat[Upper part]{\begin{tabular}{c}1\end{tabular}}\quad
  \subfloat[Lower part]{\begin{tabular}{c}2\end{tabular}}
\end{table}
\begin{figure*}[h]
  \centering
  \subfloat[Wide half]{\fbox{Wide}\label{fig:wide-a}}
  \caption{Wide figure}\label{fig:wide}
\end{figure*}
\begin{program}[h]
  \fbox{Code}
  \caption{A program}\label{prog:one}
\end{program}

Cref: \cref{fig:one}; \cref{fig:a}; \Cref{fig:b}; \cref{fig:two}; \cref{tab:one}; \cref{fig:wide}; \cref{fig:wide-a}; \cref{prog:one}; \cref{sec:res}.

Range: \cref{fig:one,fig:two,fig:wide}.

Ref: \ref{fig:one}; \ref{fig:a}; \ref{fig:two}; \ref{tab:one}; \ref{fig:wide}; \ref{prog:one}.

Autoref: \autoref{fig:one}; \autoref{tab:one}.

\emph{In italics, \cref{fig:two}.} In a footnote.\footnote{In a footnote, \cref{tab:one}.}

Subref: \subref{fig:a}; \subref{fig:b}; \subref{fig:wide-a}.

\begin{theorem}
In a theorem, \cref{fig:wide}.
\end{theorem}
\end{document}
